\addchap{Agradecimientos}

La creación de esta nueva plantilla para \LaTeX{} es el resultado del esfuerzo acumulado de numerosos compañeros de la Escuela. En primer lugar, es obligado reconocer la labor del Prof. Payán, quien, de forma desinteresada y partiendo de su experiencia en la edición de libros docentes, dio vida a la versión original de este formato; un diseño que se ha convertido en un símbolo de la imagen institucional de la ETSi. Del mismo modo, agradezco la dedicación de los profesores Juan José Murillo y Germán Madinabeitia, piezas clave en el mantenimiento de la plantilla original durante años.

Por otra parte, las portadas incluidas tienen su origen en los diseños realizados por el Prof. Fernando García García, de la Facultad de Bellas Artes, para la sección de publicaciones de nuestra Escuela. Nuestra Escuela le agradece que pusiera su arte y su trabajo, de forma gratuita, a nuestra disposición.

Finalmente, expreso mi gratitud al Profesor Antonio Estepa, tanto por su confianza para encomendarme esta labor como por las valiosas recomendaciones de redacción que ha aportado al resumen y al primer capítulo de este documento.

{\flushleft{\hfill \textit{Vicente Jesús Mayor Gallego} \\
\hfill \textit{Profesor del Departamento de Ingeniería Telemática}}}%
{\flushleft{\hfill \emph{Sevilla, 2026}}}%

\addchap{Resumen} 
La proliferación de drones o vehículos aéreos no tripulados (\glspl{uav}), tanto de uso comercial como profesional, representa una amenaza crítica para la seguridad de infraestructuras estratégicas, espacios aéreos restringidos o núcleos de población. Los sistemas actuales de detección usan distintos métodos (basados en imágenes, acústicos o láser, entre otros) para llevar a cabo la identificación de esta clase de dispositivos. De entre todos ellos, destacan las soluciones basadas en radiofrecuencia (RF) debido a que presentan características ideales para la detección temprana, como son su largo alcance y la capacidad de detección sin necesidad de tener visión directa del objetivo. Sin embargo, presentan serias limitaciones al operar en entornos urbanos o de guerra electrónica, caracterizados por un alto nivel de ruido e interferencias. Para hacer frente a estas dificultades, la tendencia actual en el desarrollo de estos sistemas se orienta hacia la creación de áreas completamente malladas con dispositivos de computación en el borde que actúen como un sistema conjunto de nodos detectores.

Gran parte de las soluciones que se muestran en la literatura actual emplean sistemas con arquitecturas de visión artificial, es decir, realizan transformaciones sobre la señal radio que llega a las antenas receptoras para facilitar su posterior procesamiento y detección. En estas operaciones, se identifican dos problemas: en primer lugar, parte de la información compleja que contiene la señal I/Q recibida se pierde al crear una representación del espectro en formato imagen; en segundo lugar, este tipo de arquitecturas emplean modelos con millones de parámetros cuyo coste computacional hace imposible su integración en nodos con capacidad de procesamiento limitada. El objetivo de este Trabajo Fin de Máster es diseñar un sistema de detección y clasificación de drones robusto, relativamente ligero y con una precisión considerable, capaz de operar bajo condiciones de degradación extrema del canal debido a bajos niveles de relación señal a ruido y la coexistencia con interferencias.
Para abordar este desafío, se propone una arquitectura de red neuronal denominada \textit{DualStream-SlidingWindow}. El sistema extrae simultáneamente las características temporales de la señal en banda base y densidad espectral de potencia. Para ello, emplea una red convolucional de valores complejos unidimensional (\textit{Complex-Values Convolutional Neural Network} 1D) combinada con mecanismos de atención para la fusión de las características obtenidas en los dominios temporales y frecuenciales. El diseño incorpora un detector de transmisiones basado en la entropía espectral de Shannon y un algoritmo de inferencia por ventana deslizante con mitigación adaptativa de ruido, permitiendo la acumulación de evidencias antes de concluir si en la ventana de señal analizada existen o no transmisiones de dron.

Las pruebas realizadas sobre el conjunto de datos \textit{NoisyUAV v2}, que cuenta con grabaciones de hasta -20 dB de relación señal a ruido (\gls{snr}) e interferencias Wi-Fi y Bluetooth, demuestran la capacidad de detección del sistema propuesto. El modelo final, con apenas 344.000 parámetros, alcanza una exactitud global del 92,20 \%, una especificidad superior al 99,8 \% y un F1-Score del 95,83 \%, logrando tasas de detección superiores al 99 \% en condiciones de $\gls{snr} \ge 0$. Adicionalmente, una validación \textit{Zero-Shot} confirmó que el sistema puede detectar la presencia de drones a partir de sus transmisiones sin que estas se incluyesen en el subconjunto de entrenamiento (79,17 \% de exhaustividad). 

Estos resultados verifican que la preservación de la fase compleja combinada con la fusión por \textit{soft-attention} ofrece una solución alternativa a los sistemas de detección basados en grandes modelos de visión artificial. El modelo propuesto supone un avance hacia el despliegue de redes de sensores \textit{Edge AI} escalables, eficientes y fiables para la detección de drones en entornos electromagnéticos hostiles.

\addchap{Abstract} 
The proliferation of Unmanned Aerial Vehicles (UAVs), for both commercial and professional use, represents a critical threat to the security of strategic infrastructures, restricted airspaces, and populated areas. Current detection systems use various methods (such as vision-based, acoustic, or laser, among others) to identify these devices. Among these, Radio Frequency (RF) based solutions stand out because they offer ideal characteristics for early detection, such as long-range capabilities and non-line-of-sight (NLOS) detection. However, they face severe limitations when operating in urban or electronic warfare environments, which are characterized by high levels of noise and interference. To address these difficulties, the current trend in the development of these systems leans towards creating fully meshed areas with edge computing devices acting as a joint system of detector nodes.

A large portion of the solutions presented in current literature employ systems with computer vision architectures; that is, they perform transformations on the radio signal arriving at the receiving antennas to facilitate its subsequent processing and detection. Two problems are identified in these operations: first, part of the complex information contained in the received I/Q signal is lost when creating an image-format representation of the spectrum; second, this type of architecture employs models with millions of parameters that are computationally unfeasible to embed in nodes with limited processing capacity. The objective of this Master's Thesis is to design a robust, relatively lightweight, and highly accurate drone detection and classification system, capable of operating under extreme channel degradation conditions due to low signal-to-noise ratios and the coexistence of interference.

To address this challenge, a neural network architecture named \textit{DualStream-SlidingWindow} is proposed. The system simultaneously extracts the temporal features of the baseband signal (I/Q) and its spectral morphology (Power Spectral Density). To achieve this, it employs a 1D Complex-Valued Convolutional Neural Network (CV-CNN) combined with attention mechanisms to fuse the features obtained from both domains. Operationally, the design incorporates a transmission detector based on Shannon's spectral entropy and a sliding window inference algorithm with adaptive noise mitigation, allowing the accumulation of evidence before concluding whether a drone transmission exists in the analyzed signal.

The experiments, conducted on the \textit{NoisyUAV v2} dataset (which includes recordings with \glspl{snr} down to -20 dB and Wi-Fi/Bluetooth interference), demonstrate the detection capability of the proposed system. The final model, with only 344,000 parameters, achieves an overall accuracy of 92.20\% and an F1-Score of 95.83\%, reaching detection rates over 99\% under favorable \gls{snr} conditions. Even in extreme noise regimes, the system maintains high specificity ($>99.8\%$) under a configuration termed \textit{Maximum Security}. Additionally, a \textit{Zero-Shot} validation confirmed the network's ability to detect drone transmissions on which it had never been trained (79.17\% \textit{Recall}). 

These results verify that the preservation of complex phase combined with \textit{soft-attention} fusion offers an alternative solution to detection systems based on large computer vision models. The proposed system represents a significant step towards the deployment of scalable, efficient, and reliable \textit{Edge AI} sensor networks for \gls{uav} detection in hostile electromagnetic environments.
